\documentclass[11pt]{article}

\usepackage[margin=1in]{geometry}
\usepackage{hyperref}
\usepackage{enumitem}
\usepackage{xcolor}
\usepackage{titlesec}
\usepackage{array}
\usepackage{tcolorbox}
\usepackage{parskip}

\definecolor{primary}{HTML}{6366F1}
\definecolor{secondary}{HTML}{06B6D4}
\definecolor{dark}{HTML}{111827}
\definecolor{light}{HTML}{F3F4F6}

\hypersetup{
    colorlinks=true,
    linkcolor=primary,
    urlcolor=primary
}

\titleformat{\section}
{\Large\bfseries\color{primary}}
{}
{0em}
{}

\titleformat{\subsection}
{\bfseries\color{dark}}
{}
{0em}
{}

\setlist[itemize]{leftmargin=1.5em,itemsep=0.45em}

\begin{document}

\begin{titlepage}
    \centering
    \vspace*{2cm}

    {\Huge\bfseries PRD Autopilot\par}
    \vspace{0.5cm}
    {\Large\color{primary} AI-Powered Founder Platform\par}

    \vspace{1cm}

    \begin{tcolorbox}[
        colback=light,
        colframe=primary,
        width=0.85\textwidth,
        arc=4mm,
        boxrule=0.8pt
    ]
    \centering
    \textbf{Turn ideas into engineering-ready PRDs, validation reports, and investor pitch decks in minutes.}
    \end{tcolorbox}

    \vspace{1.5cm}

    {\large MTP World Product Day 2026\par}
    \vspace{0.3cm}
    {\large June 2026\par}

    \vfill

    {\large\bfseries Team Members\par}
    \vspace{0.3cm}
    Erfan \quad $\bullet$ \quad Elijah \quad $\bullet$ \quad Esmeralda

    \vspace{1cm}
\end{titlepage}

\section{Short Description}

PRD Autopilot is an AI powered platform that helps users validate ideas, generate Product Requirements Documents (PRDs), analyze product quality, create investor ready pitch decks, and collaborate with AI throughout the product development process.

\section{Project Highlights}

\begin{tcolorbox}[
    colback=light,
    colframe=secondary,
    arc=3mm,
    boxrule=0.8pt
]
\centering
\textbf{Idea Validation}
$\rightarrow$
\textbf{PRD Generation}
$\rightarrow$
\textbf{Health Scoring}
$\rightarrow$
\textbf{Backlog Creation}
$\rightarrow$
\textbf{Pitch Decks}
$\rightarrow$
\textbf{AI Chat}
\end{tcolorbox}

\begin{itemize}
    \item Multi agent AI PRD generation
    \item AI powered idea validation reports
    \item PRD health scoring and gap analysis
    \item Ticket and backlog generation
    \item Investor ready pitch deck generation
    \item AI product advisor chat
    \item PDF, Markdown, and Notion export
\end{itemize}

\section{Inspiration}

Many founders, students, product managers, and startup teams have great ideas but struggle to turn them into structured product plans. Researching competitors, validating markets, writing PRDs, and creating pitch decks can take days or weeks.

We wanted to build a platform that acts like an AI co-founder, helping users move from a rough idea to a validated, investor ready product strategy in minutes.

\section{What It Does}

PRD Autopilot helps users:

\begin{itemize}
    \item Validate startup and product ideas using AI powered market research.
    \item Generate complete Product Requirements Documents through an AI interview process.
    \item Score PRD quality and identify gaps.
    \item Create structured backlogs and product tickets.
    \item Generate investor ready pitch decks.
    \item Chat with an AI product advisor for guidance and iteration.
    \item Export work to PDF, Markdown, and Notion.
\end{itemize}

The platform transforms a rough concept into a complete product planning package.

\section{How We Built It}

PRD Autopilot was built using Next.js, TypeScript, Tailwind CSS, Supabase, Groq, Vercel, and the Notion API.

The platform combines multiple AI powered workflows including idea validation, PRD generation, PRD health scoring, backlog generation, pitch deck creation, and AI assisted product guidance.

User authentication, dashboards, generated documents, validation reports, pitch decks, and project data are handled through Supabase. AI generated content is powered through Groq and specialized prompt pipelines designed for different parts of the founder workflow.

\section{Challenges We Ran Into}

One of the biggest challenges was generating structured AI outputs that remained consistent across many different product categories and industries. We had to make sure the AI could produce useful PRDs, validation reports, pitch decks, personas, acceptance criteria, and success metrics without returning broken or overly generic content.

We also ran into challenges with authentication, session persistence, protected routes, AI rate limits, PDF and Notion exports, and making sure generated content stayed connected to the correct user account.

Another challenge was testing the platform across many different ideas. Some outputs were strong, while others revealed issues such as incomplete validation sections, generic edge cases, or titles being pulled from interview answers instead of generated as proper product names.

\section{Accomplishments We Are Proud Of}

\begin{itemize}
    \item Built a complete AI powered founder workflow from idea validation to pitch deck generation.
    \item Created a PRD generation system that turns a rough idea and five follow up questions into a structured product document.
    \item Added PRD health scoring and gap analysis.
    \item Built AI powered idea validation reports with market, competitor, risk, and SWOT analysis.
    \item Generated investor ready pitch decks from ideas, PRDs, or validation reports.
    \item Added export support for PDF, Markdown, and Notion.
    \item Built authentication, protected routes, user dashboards, and saved project history.
    \item Deployed the app to Vercel and tested it across desktop and mobile.
\end{itemize}

\section{What We Learned}

We learned how important prompt engineering, structured AI responses, and real user testing are when building AI powered tools. Even when the app worked technically, testing showed us where AI output could be improved, where UX could be clearer, and where edge cases could break the experience.

We also learned how to combine multiple AI workflows into one platform. Instead of only generating a document, PRD Autopilot supports a larger product development flow: validating an idea, generating a PRD, scoring the PRD, creating tickets, generating a pitch deck, and chatting with AI for iteration.

\section{What's Next}

\begin{itemize}
    \item Team collaboration and shared workspaces.
    \item More advanced PRD editing workflows.
    \item Google OAuth support.
    \item Integrations with tools like Jira, Linear, Trello, and Slack.
    \item More customizable pitch decks.
    \item Stronger validation reports with deeper market intelligence.
    \item Better AI coaching for product strategy and founder decision making.
    \item Enterprise sharing, permissions, and stakeholder feedback workflows.
\end{itemize}

\section{Tech Stack}

\begin{tcolorbox}[
    colback=light,
    colframe=primary,
    arc=3mm,
    boxrule=0.8pt
]
\begin{tabular}{>{\bfseries}p{4cm} p{10cm}}
Frontend & Next.js 14, TypeScript, Tailwind CSS \\
Backend and Database & Supabase \\
AI & Groq, Llama 3.3 70B \\
Deployment & Vercel \\
Integrations & Notion API \\
\end{tabular}
\end{tcolorbox}

\section{Links}

\begin{itemize}
    \item \textbf{GitHub Repository:} \href{https://github.com/erfwn81/PRD-Autopilot}{https://github.com/erfwn81/PRD-Autopilot}
    \item \textbf{Live Demo:} \href{https://prd-autopilot.vercel.app}{https://prd-autopilot.vercel.app}
    \item \textbf{Demo Video:} To be added after recording.
\end{itemize}

\section{Team Members}

\begin{center}
\begin{tabular}{|>{\bfseries}p{3cm}|p{10cm}|}
\hline
Name & Responsibilities \\
\hline
Erfan & Backend, AI systems, architecture, deployment \\
\hline
Elijah & Frontend, dashboard, UI/UX, mobile polish \\
\hline
Esmeralda & Authentication, QA testing, demo video, Devpost submission \\
\hline
\end{tabular}
\end{center}

\end{document}